\section*{Results}

\subsection*{What the labels cover}
The external intervention cohort comprised \ExtNOne{} informative studies
(\ExtNOnePatients{} patients). Label availability differed sharply between
sites (\cref{fig:labels}; \stab{stab:labels}). With impression labels,
\NoLabSrcImpression{} of source and \NoLabExtImpression{} of external studies
had no positive or negative label for any target. Similar aggregate positive
rates (\PosSrcImpression{} and \PosExtImpression{}) concealed differences in
composition. External labels carried far more uncertain mentions, and case mix
differed: counting unmentioned findings as negative, effusion prevalence was
\SrcEffUnmNeg{} at the source and \ExtEffUnmNeg{} externally. With findings
labels the asymmetry reversed. Only \NoLabSrcFindings{} of the restricted
source cohort lacked a labelled target, against \NoLabExtFindings{} of the
external cohort. Where both sections labelled a cell, impression and findings
labels agreed in \AgreeMinSrc{}--\AgreeMaxExt{} of cells. The two endpoints
differ mainly in which cells carry a label.

\subsection*{Registered hypotheses after correction}
\Cref{tab:hyp} reports every originally computed contrast under both
estimators. The original estimator confirmed H1: \HOneMThreeOrig{} for late
fusion (M3) and \HOneMFourOrig{} for feature fusion (M4). Over evaluable studies,
the gaps were \HOneMThreeEval{} and \HOneMFourEval{}, so H1 was not confirmed.
Both intervals lie inside $\pm$0.02, conditional on the calibration sample and
training run. H2 was not confirmed under either estimator, and neither was the
registered H3, which equals H2. The post hoc transfer-gap H3 was negative for
both models; corrected, its magnitude fell below 0.02, so the previously
reported ``reversal at materiality'' does not survive.

\begin{table}[H]
\centering\footnotesize
\caption{Contrasts under the original and corrected (evaluable-study)
estimators, with verdicts (V) under the registered rule. C, confirmed; S, same
direction but below materiality; I, interval includes zero; O, opposite
direction (below materiality, or none registered); O*, opposite direction with
magnitude $\ge$0.02. H3 brackets show the verdict with the 0.02 materiality
applied by the original code. Values are external minus source unless
labelled C2 $-$ C1.}
\label{tab:hyp}
\begin{adjustbox}{max width=\linewidth}
\begin{tabular}{lllrlrl}
\toprule
H & Model & Contrast & \multicolumn{2}{c}{Original estimator} & \multicolumn{2}{c}{Corrected (evaluable-study)}\\
\cmidrule(lr){4-5}\cmidrule(lr){6-7}
 & & & Estimate (95\% CI) & V & Estimate (95\% CI) & V\\
\midrule
H1 & M3 & ext $-$ src, C0 & +0.0469 (+0.0408, +0.0531) & C & $-$0.0042 ($-$0.0145, +0.0061) & I\\
H1 & M4 & ext $-$ src, C0 & +0.0384 (+0.0336, +0.0430) & C & $-$0.0018 ($-$0.0108, +0.0064) & I\\
H2 & M3 & interaction, C1 & +0.0024 ($-$0.0027, +0.0076) & I & +0.0048 ($-$0.0038, +0.0134) & I\\
H2 & M4 & interaction, C1 & +0.0096 (+0.0048, +0.0144) & S & +0.0089 (+0.0002, +0.0173) & S\\
\midrule\multicolumn{7}{l}{\emph{H3 as coded after inspection (registered H3 $\equiv$ H2-C1)}}\\
H3 & M3 & gap $-$ gap(M1) & $-$0.0265 ($-$0.0321, $-$0.0212) & O [O*] & $-$0.0162 ($-$0.0256, $-$0.0070) & O [O]\\
H3 & M4 & gap $-$ gap(M1) & $-$0.0350 ($-$0.0408, $-$0.0292) & O [O*] & $-$0.0138 ($-$0.0243, $-$0.0029) & O [O]\\
\midrule\multicolumn{7}{l}{\emph{Secondary: H4 (direction only, no materiality)}}\\
H4 & M3 & C2 $-$ C1, source & +0.0475 (+0.0422, +0.0524) & C & +0.0783 (+0.0691, +0.0865) & C\\
H4 & M4 & C2 $-$ C1, source & +0.0173 (+0.0125, +0.0217) & C & +0.0301 (+0.0215, +0.0383) & C\\
H4 & M3 & C2 $-$ C1, external & +0.0231 (+0.0211, +0.0250) & C & +0.0263 (+0.0240, +0.0286) & C\\
H4 & M4 & C2 $-$ C1, external & $-$0.0060 ($-$0.0080, $-$0.0041) & O & $-$0.0077 ($-$0.0101, $-$0.0054) & O\\
\midrule\multicolumn{7}{l}{\emph{Not registered (computed by the original analysis code; descriptive)}}\\
H2 & M3 & interaction, C2 & $-$0.0219 ($-$0.0279, $-$0.0160) & O* & $-$0.0472 ($-$0.0572, $-$0.0374) & O*\\
H2 & M4 & interaction, C2 & $-$0.0137 ($-$0.0185, $-$0.0087) & O & $-$0.0289 ($-$0.0377, $-$0.0199) & O*\\
H1 & M1 & ext $-$ src, C0 & +0.0734 (+0.0668, +0.0802) & C & +0.0120 ($-$0.0002, +0.0244) & I\\
H1 & M2 & ext $-$ src, C0 & +0.0763 (+0.0687, +0.0838) & C & $-$0.0173 ($-$0.0292, $-$0.0054) & O\\
\bottomrule
\end{tabular}

\end{adjustbox}
\end{table}

\subsection*{The cross-site contrast depends on analysis choices}
With predictions and frozen thresholds fixed, label handling alone moved the
M4 transfer gap from \SensUncPosHOneMFour{} (uncertain counted positive) to
\SensUnmNegHOneMFour{} (unmentioned counted negative). For M3 the range was
\SensTTenHOneMThreePt{} to \SensUnmNegHOneMThreePt{} (\cref{fig:spec};
\stab{stab:sens}). A second training seed moved the corrected gaps to
\SensSeedHOneMThree{} for M3 and \SensSeedHOneMFour{} for M4. Aggregation, label
source and view stratum each moved at least one model's estimate by more than
its interval width. M4's gap was \ViewPAMFour{} in PA-only studies and
\ViewAPMFour{} in AP-only studies, and AP films made up \APShareSrc{} of source
and \APShareExt{} of external studies. Pathology-specific gaps for M4 were
heterogeneous under the primary labels: \PathMFourCard{} for cardiomegaly and
\PathMFourAtel{} for atelectasis (\ssec{ssec:path}).

\subsection*{Controlled label-source decomposition}
The original findings-label sensitivity changed four things at once: labels,
source cohort, thresholds and C2 donors. Its reported reversal of H1 reproduces
exactly (\cref{tab:ls}). Changing one factor at a time showed that, under the
original estimator, swapping labels on the same studies (A0$\to$A1) and
restricting the source cohort (A1$\to$A2) produced most of the reversal. The
restriction matters only because it removes label-less source studies that the
defect counted as correct; under the corrected estimator it has exactly no
effect. Under the corrected estimator, the findings-label result does not
reverse the primary: \LSAThreeMThreeEval{} for M3 and \LSAThreeMFourEval{} for
M4. On cells labelled under both sources, the two label sets gave nearly the
same contrasts (steps A1c). Re-selecting thresholds on findings labels also
changed external decisions in \SOneCellsChanged{} cells, so the earlier
argument that ``only the labels differ'' did not hold.

\begin{table}[H]
\centering\footnotesize
\caption{Controlled label-source decomposition of H1 (point estimates;
intervals in the supplement). Evaluable: share of studies with at least one
labelled target.}
\label{tab:ls}
\begin{adjustbox}{max width=\linewidth}
\begin{tabular}{llrrrrrr}
\toprule
 & & \multicolumn{2}{c}{Evaluable studies} & \multicolumn{2}{c}{H1 M3} & \multicolumn{2}{c}{H1 M4}\\
\cmidrule(lr){3-4}\cmidrule(lr){5-6}\cmidrule(lr){7-8}
Step & Change from A0 & Src & Ext & Original & Corrected & Original & Corrected\\
\midrule
A0 & primary (impression labels, primary cohort, impression thresholds) & 53\% & 81\% & +0.0469 & $-$0.0042 & +0.0384 & $-$0.0018\\
A1 & findings labels only & 59\% & 24\% & $-$0.0613 & +0.0090 & $-$0.0839 & $-$0.0389\\
A1c-I & common labelled cells, impression labels & 20\% & 17\% & $-$0.0059 & $-$0.0042 & $-$0.0065 & $-$0.0212\\
A1c-F & common labelled cells, findings labels & 20\% & 17\% & $-$0.0069 & $-$0.0090 & $-$0.0077 & $-$0.0271\\
A2x & findings cohort only & 37\% & 81\% & +0.0809 & +0.0012 & +0.0664 & +0.0126\\
A2 & findings cohort + labels & 94\% & 24\% & $-$0.1359 & +0.0090 & $-$0.1610 & $-$0.0389\\
A3 & + findings thresholds and cutoffs (= original S1) & 94\% & 24\% & $-$0.1324 & +0.0285 & $-$0.1079 & +0.0048\\
B2 & findings labels + thresholds, primary cohort & 59\% & 24\% & $-$0.0579 & +0.0285 & $-$0.0524 & +0.0048\\
\bottomrule
\end{tabular}

\end{adjustbox}
\end{table}

\subsection*{Operating points, coverage and alternative scores}
Frozen thresholds transferred a decision rule, not an operating point
(\cref{fig:operating}; \stab{stab:operating}). For M4 among accepted studies,
cardiomegaly specificity fell from \MFourCardSpecSrc{} to \MFourCardSpecExt{}
and positive calls rose from \MFourCardPosSrc{} to \MFourCardPosExt{}.
Specificity also fell for effusion and edema. This direction held when
unmentioned cells counted as negative. Negative labels are explicit mentions,
so these specificities carry verification bias. With context intact, realised
coverage stayed within \CovDriftMaxCZero{} of target at both sites, and within
\SeedCovDriftMax{} under the second seed. Under intervention it did not. M2
accepted every study under C1, because a constant placeholder makes its output
constant. M3 coverage under C2 fell to \MThreeCTwoCovSrc{} at the source and
\MThreeCTwoCovExt{} externally. The frozen score
ranked study-level correctness weakly, with failure-detection AUROC
\FailAUROCFrozenMin{}--\FailAUROCFrozenMax{} (\cref{fig:rc}). A decision-aware
expected-loss score fitted on the source calibration tier ranked errors better
(\FailAUROCELMin{}--\FailAUROCELMax{}). Its coverage transported worse, reaching
\ELMFourExtCov{} externally for M4 (\stab{stab:comparators}).

\subsection*{Context interventions within each site}
For M3, misaligned context gave higher error than absent context at the source
under every specification, C2 permutation and both training seeds
(\HFourMThreesourceEval{}; \cref{fig:h4}). Externally it did so under every
specification and C2 permutation for the first seed (\HFourMThreeexternalEval{}),
but not under the second seed (\RepTwoHFourMThreeexternaleval{}). For M4 the
source effect was positive except under the second seed; externally its sign
depended on label handling. Interventions acted mainly through changed
predictions. Even so, accepted sets under C0 and intervention overlapped only
partially (Jaccard \JaccardMin{}--\JaccardMax{}; \stab{stab:context}).

\subsection*{Uncertainty beyond the evaluation bootstrap}
Resampling calibration patients and re-applying the frozen selection rule
moved corrected H1 estimates with standard deviations
\CalRatioMin{}--\CalRatioMax{} times the evaluation-bootstrap standard
deviations (\cref{tab:unc}). For M4 the calibration standard deviation was
\CalSDMFour{} against \EvalSDMFour{} from the evaluation bootstrap, and
per-pathology threshold standard deviations reached \ThrSDMaxMFour{}. The
reported intervals therefore understate total uncertainty for cross-site
contrasts.

\begin{table}[H]
\centering\footnotesize
\caption{H1 variability from resampling evaluation patients (bootstrap SD,
approximated from the 95\% interval) and from resampling calibration patients
with the frozen selection rule re-applied (200 replicates).}
\label{tab:unc}
\begin{adjustbox}{max width=\linewidth}
\begin{tabular}{llrrrr}
\toprule
Model & Estimator & H1 & SD (evaluation bootstrap) & SD (calibration resample) & Calibration 2.5--97.5\%\\
\midrule
M1 & original & +0.0734 & 0.0034 & 0.0044 & (+0.0635, +0.0845)\\
 & evaluable-study & +0.0120 & 0.0063 & 0.0092 & ($-$0.0003, +0.0388)\\
M2 & original & +0.0763 & 0.0038 & 0.0067 & (+0.0740, +0.0984)\\
 & evaluable-study & $-$0.0173 & 0.0061 & 0.0205 & ($-$0.0254, +0.0408)\\
M3 & original & +0.0469 & 0.0031 & 0.0050 & (+0.0352, +0.0554)\\
 & evaluable-study & $-$0.0042 & 0.0053 & 0.0097 & ($-$0.0271, +0.0141)\\
M4 & original & +0.0384 & 0.0024 & 0.0032 & (+0.0316, +0.0440)\\
 & evaluable-study & $-$0.0018 & 0.0044 & 0.0102 & ($-$0.0286, +0.0102)\\
\bottomrule
\end{tabular}

\end{adjustbox}
\end{table}

\subsection*{Integrity checks}
The previously reported studies ``retaining their own context'' under C2
(\CTwoIdentSrc{} source, \CTwoIdentExt{} external) all received identical text
from a different patient; no recipient received its own patient's context.
Exploratory external demographic strata are in \stab{stab:demo}.
